% GENERATED FILE. Edit canonical lessons, then run npm run generate:books.
% canonical-source-hash: fnv1a64:cd3773ca88594c00
% canonical-lessons: MW-C156-mhari, MW-C156-maay-mathai, MW-C156-su, MW-R156-mine-in-from

\chapter{Mine, In and From}
\label{ch:mw-a2-mine-in-and-from}
\hlchaptermodality{156}

\begin{chapteropening}
\textbf{By the end of this chapter:} \emph{I can say whose a thing is, where it is, and where someone comes from.}

Complete MW-R156-mine-in-from: make my and your agree, put things in and on places, and say where someone comes from.
\end{chapteropening}

\section[\texorpdfstring{\mw{म्हारो} · \mw{म्हारी} (mhāro · mhārī)}{म्हारो · म्हारी (mhāro · mhārī)}]{\mw{म्हारो}, \mw{म्हारी} — my, agreeing with what is mine}

\label{lesson:MW-C156-mhari}



\begin{quote}
Say \emph{my name is Ram}, then \emph{this is my house}. Then say \emph{a boy} and \emph{a girl}: \textbf{\mw{छोरो}}, \textbf{\mw{छोरी}}.
\end{quote}

\begin{grammarlens}[title={Grammar Lens: the thing decides}]
\textbf{\mw{म्हारो}} ends like \textbf{\mw{छोरो}}, and it changes the way \textbf{\mw{छोरो}} turns into \textbf{\mw{छोरी}}: it follows the thing that is yours.

\begin{quote}
\textbf{\mw{यो} \mw{म्हारो} \mw{घर} \mw{है।}}
\end{quote}

\begin{quote}
\emph{yo mhāro ghar hai.} — This is my house.
\end{quote}

\begin{quote}
\textbf{\mw{म्हारी} \mw{गाय} \mw{अठै} \mw{है।}}
\end{quote}

\begin{quote}
\emph{mhārī gāy aṭhai hai.} — My cow is here.
\end{quote}

\textbf{\mw{घर}} is masculine, so \textbf{\mw{म्हारो}}; \textbf{\mw{गाय}} is feminine, so \textbf{\mw{म्हारी}}. A woman says \textbf{\mw{म्हारो} \mw{घर}} too: the ending follows the house, not the owner.
\end{grammarlens}

\subsection*{Your turn}
\textbf{\mw{थारो}}, \emph{your}, does the same: \textbf{\mw{थारो} \mw{दूध}}, \textbf{\mw{थारी} \mw{रोटी}}.

Say these aloud:

\begin{itemize}
\raggedright
  \item \textbf{\mw{म्हारी} \mw{रोटी} \mw{गरम} \mw{है।}} — \emph{mhārī roṭī garam hai.} — My roti is hot.
  \item \textbf{\mw{थारो} \mw{दूध} \mw{अठै} \mw{है।}} — \emph{thāro dūdh aṭhai hai.} — Your milk is here.
  \item \textbf{\mw{म्हारो} \mw{भाई} \mw{बठै} \mw{है।}} — \emph{mhāro bhāī baṭhai hai.} — My brother is there.
\end{itemize}

\begin{tcolorbox}[breakable,colback=teal!4,colframe=teal!35!black,title={Before you move on}]
\emph{My house}? (\textbf{\mw{म्हारो} \mw{घर}}.) \emph{My cow}? (\textbf{\mw{म्हारी} \mw{गाय}}.) Who decides the ending? (The thing that is yours.)
\end{tcolorbox}
\section[\texorpdfstring{\mw{मांय} · \mw{माथै} (mā̃y · māthai)}{मांय · माथै (mā̃y · māthai)}]{\mw{रसोई} \mw{मांय}, \mw{मेज} \mw{माथै} — in and on}

\label{lesson:MW-C156-maay-mathai}



\begin{quote}
Say \emph{inside} and \emph{on top}: \textbf{\mw{मांय}}, \textbf{\mw{माथै}}. Then \emph{my cow is here}.
\end{quote}

\begin{grammarlens}[title={Grammar Lens: a noun in front}]
Put a noun in front of \textbf{\mw{मांय}} and \textbf{\mw{माथै}}, and they say \emph{in} and \emph{on}:

\begin{quote}
\textbf{\mw{दूध} \mw{रसोई} \mw{मांय} \mw{है।}}
\end{quote}

\begin{quote}
\emph{dūdh rasoī mā̃y hai.} — The milk is in the kitchen.
\end{quote}

\begin{quote}
\textbf{\mw{किताब} \mw{मेज} \mw{माथै} \mw{है।}}
\end{quote}

\begin{quote}
\emph{kitāb mej māthai hai.} — The book is on the table.
\end{quote}

\textbf{\mw{कनै}}, \emph{near}, works the same: \textbf{\mw{घर} \mw{कनै}}, near the house. Old Marwari already had \emph{kanai}, \textquotedblleft{}near\textquotedblright{}, and \emph{māhī}, \textquotedblleft{}in, into\textquotedblright{} (CDIAL 2830, 9804).
\end{grammarlens}

\subsection*{Your turn}
Say these aloud:

\begin{itemize}
\raggedright
  \item \textbf{\mw{दुकान} \mw{बाजार} \mw{मांय} \mw{है।}} — \emph{dukān bājār mā̃y hai.} — The shop is in the market.
  \item \textbf{\mw{चाय} \mw{मेज} \mw{माथै} \mw{है।}} — \emph{chāy mej māthai hai.} — The tea is on the table.
  \item \textbf{\mw{वो} \mw{घर} \mw{कनै} \mw{है।}} — \emph{vo ghar kanai hai.} — He is near the house.
\end{itemize}

\begin{tcolorbox}[breakable,colback=teal!4,colframe=teal!35!black,title={Before you move on}]
\emph{In the kitchen}? (\textbf{\mw{रसोई} \mw{मांय}}.) \emph{On the table}? (\textbf{\mw{मेज} \mw{माथै}}.)
\end{tcolorbox}
\section[\texorpdfstring{\mw{सूं} (sū̃)}{सूं (sū̃)}]{\mw{बाजार} \mw{सूं} — from}

\label{lesson:MW-C156-su}



\begin{quote}
Say \emph{the shop is in the market}. Then \emph{he comes every day}, and the question word \emph{where}, \textbf{\mw{कठै}}.
\end{quote}

\begin{grammarlens}[title={Grammar Lens: \mw{सूं}}]
\textbf{\mw{सूं}} after a place says where someone comes \emph{from}:

\begin{quote}
\textbf{\mw{वो} \mw{बाजार} \mw{सूं} \mw{आवै} \mw{है।}}
\end{quote}

\begin{quote}
\emph{vo bājār sū̃ āvai hai.} — He comes from the market.
\end{quote}

\begin{quote}
\textbf{\mw{हूं} \mw{दुकान} \mw{सूं} \mw{आऊं} \mw{हूं।}}
\end{quote}

\begin{quote}
\emph{hū̃ dukān sū̃ āū̃ hū̃.} — I come from the shop.
\end{quote}

Old Marwari wrote it \emph{-su}: \emph{talā-su}, \textquotedblleft{}from below\textquotedblright{} (CDIAL 5731).
\end{grammarlens}

\subsection*{Your turn}
Ask with \textbf{\mw{कठै}} in front of it:

\begin{quote}
\textbf{\mw{आप} \mw{कठै} \mw{सूं} \mw{आवो} \mw{हो}?}
\end{quote}

\begin{quote}
\emph{āp kaṭhai sū̃ āvo ho?} — Where do you come from?
\end{quote}

Say these aloud:

\begin{itemize}
\raggedright
  \item \textbf{\mw{आप} \mw{कठै} \mw{सूं} \mw{आवो} \mw{हो}?} — \emph{āp kaṭhai sū̃ āvo ho?} — Where do you come from?
  \item \textbf{\mw{हूं} \mw{बाजार} \mw{सूं} \mw{आऊं} \mw{हूं।}} — \emph{hū̃ bājār sū̃ āū̃ hū̃.} — I come from the market.
\end{itemize}

\begin{tcolorbox}[breakable,colback=teal!4,colframe=teal!35!black,title={Before you move on}]
\emph{From the market}? (\textbf{\mw{बाजार} \mw{सूं}}.) \emph{Where from}? (\textbf{\mw{कठै} \mw{सूं}}.)
\end{tcolorbox}
\section[(dialogue)]{Mine, in and from — a review}

\label{lesson:MW-R156-mine-in-from}



\begin{quote}
Say \emph{my cow is here}.
\end{quote}

\subsection*{Your turn}
Say these aloud:

\begin{itemize}
\raggedright
  \item \emph{my house}; \emph{my cow}; \emph{your roti}; \emph{your milk}
  \item \emph{the milk is in the kitchen}; \emph{the book is on the table}
  \item \emph{the shop is near the house}
  \item \emph{where do you come from?}; \emph{I come from the market}
  \item \emph{he comes from the shop every day}
\end{itemize}

\begin{tcolorbox}[breakable,colback=teal!4,colframe=teal!35!black,title={Before you move on}]
\emph{My cow}? (\textbf{\mw{म्हारी} \mw{गाय}}.) \emph{In the kitchen}? (\textbf{\mw{रसोई} \mw{मांय}}.) \emph{From the shop}? (\textbf{\mw{दुकान} \mw{सूं}}.)
\end{tcolorbox}
